\chapter{DEPLOYMENT}

Deployment is the final phase of the development ecosystem where the web application is made available over the internet for end-users. The Alumni Management System deployment strategy ensures high availability, security, and scalability.

\section{Hosting Environment}
For production deployment, a Linux-based Virtual Private Server (VPS), such as an Ubuntu instance on AWS EC2 or DigitalOcean, is highly recommended. 
\begin{itemize}
    \item \textbf{Web Server:} Nginx or Apache serves as the primary web server to handle incoming HTTP requests and serve static assets compiled by Vite.
    \item \textbf{Database:} A managed MySQL database service or a locally installed MySQL database is used to store persistent data securely.
    \item \textbf{PHP Processing:} PHP-FPM manages back-end code execution efficiently.
\end{itemize}

\section{Deployment Process}
The deployment pipeline follows modern DevOps practices:
\begin{enumerate}
    \item \textbf{Version Control Integration:} Code is pulled directly from the main repository (e.g., GitHub or GitLab) onto the production server.
    \item \textbf{Dependency Management:} \texttt{composer install --no-dev} and \texttt{npm install} are run to install necessary PHP and JavaScript packages respectively, omitting development-only dependencies.
    \item \textbf{Asset Compilation:} \texttt{npm run build} compiles and minifies Tailwind CSS and JavaScript assets using Vite.
    \item \textbf{Environment Configuration:} The \texttt{.env} file is securely configured with production credentials, ensuring debug mode is set to \texttt{false}.
    \item \textbf{Database Migration:} \texttt{php artisan migrate} prepares the database schema, and necessary initial seeding (e.g., Admin user creation) is completed.
    \item \textbf{Cache Optimization:} Configuration, routing, and view layers are cached using artisan commands (\texttt{php artisan optimize}) to significantly improve response times.
\end{enumerate}

\section{Folder Structure}
The project follows standard Laravel MVC conventions, enabling a clean separation of concerns.

\subsection{Key Directories}
\begin{itemize}
    \item \textbf{app/}: Contains the core code of the application.
    \begin{itemize}
        \item \textbf{Http/Controllers/}: Holds the business logic for responding to user requests.
        \item \textbf{Http/Middleware/}: Access filters (e.g., Auth, Role Management, Profile Completion).
        \item \textbf{Models/}: Eloquent classes defining database interactions.
    \end{itemize}
    \item \textbf{bootstrap/}: App initialization scripts and autoloading.
    \item \textbf{config/}: Holds all central configuration files (database, cache, mail).
    \item \textbf{database/}: Houses migrations (schema builders), model factories, and seeders.
    \item \textbf{public/}: The \texttt{index.php} file serves as the singular entry point. Stores compiled assets via Vite.
    \item \textbf{resources/}: Contains raw front-end assets.
    \begin{itemize}
        \item \textbf{views/}: The Blade HTML templates.
        \item \textbf{css/ \& js/}: Raw Tailwind CSS and JavaScript logic before Vite compilation.
    \end{itemize}
    \item \textbf{routes/}: Defines all HTTP web and API endpoints. 
    \item \textbf{storage/}: Stores compiled Blade templates, logs, file sessions, and user-uploaded media (e.g., profile pictures, post images).
    \item \textbf{tests/}: Contains automated test suites (Pest/PHPUnit).
\end{itemize}
